\subsection{Equations}

Each equation can be expressed in either integral or differential form.
I'll present them in their differential form. We'll use 3D Cartesian
coordinates \((x,y,z)\) and time \(t\) to describe the functions.

\subsubsection{Faraday's Law of Induction}
Faraday's Law of Induction states that a time-varying magnetic function
induces an electric function. It is expressed as:
\begin{align*}
  -\nabla \times \mathbf{E(x,y,z,t)} = \frac{\partial \mathbf{B(x,y,z,t)}}{\partial t} + \mathbf{J_m(x,y,z,t)} \tag{1}
\end{align*}
where
\begin{itemize}
  \item \( \nabla \):
    \begin{align*}
      \nabla = \left( \frac{\partial}{\partial x}, \frac{\partial}{\partial y}, \frac{\partial}{\partial z} \right)
    \end{align*}
  \item \( \mathbf{E(x,y,z,t)} \): electric function. \( \frac{V}{m} \)
  \item \( \mathbf{B(x,y,z,t)} \): magnetic flux density function.
    \( \frac{Wb}{m^2} \)
  \item \( \mathbf{J_m(x,y,z,t)} \): magnetic current density function.
    \( \frac{A}{m^2} \). In nature, \( \mathbf{J_m(x,y,z,t)} = 0 \) since
    magnetic monopoles that is, charges haven't been observed.
\end{itemize}
Let's have some fun with this equation.
\begin{align*}
  \text{[Divergence of both sides]} \\
  -\nabla \cdot (\nabla \times \mathbf{E(x,y,z,t)}) &= \nabla \cdot \left(\frac{\partial \mathbf{B(x,y,z,t)}}{\partial t} + \mathbf{J_m(x,y,z,t)}\right) \\
  -\nabla \cdot (\nabla \times \mathbf{E(x,y,z,t)}) &= \nabla \cdot \frac{\partial \mathbf{B(x,y,z,t)}}{\partial t} + \nabla \cdot \mathbf{J_m(x,y,z,t)} \\
  -\nabla \cdot (\nabla \times \mathbf{E(x,y,z,t)}) &= \frac{\partial}{\partial t} (\nabla \cdot \mathbf{B(x,y,z,t)}) + \nabla \cdot \mathbf{J_m(x,y,z,t)} \\
  0 &= \frac{\partial}{\partial t} (\nabla \cdot \mathbf{B(x,y,z,t)}) + \nabla \cdot \mathbf{J_m(x,y,z,t)} \\
  - \frac{\partial}{\partial t} (\nabla \cdot \mathbf{B(x,y,z,t)}) &= \nabla \cdot \mathbf{J_m(x,y,z,t)} \\
  \text{[Integrating both sides over time]} \\
  \int_{t_0}^{t_1} - \frac{\partial}{\partial t} (\nabla \cdot \mathbf{B(x,y,z,t)}) \, dt &= \int_{t_0}^{t_1} \nabla \cdot \mathbf{J_m(x,y,z,t)} \, dt \\
  - \left[\nabla \cdot \mathbf{B(x,y,z,t)}\right]_{t_0}^{t_1} &= \int_{t_0}^{t_1} \nabla \cdot \mathbf{J_m(x,y,z,t)} \, dt \\
  - \left[\nabla \cdot \mathbf{B(x,y,z,t_1)} - \nabla \cdot \mathbf{B(x,y,z,t_0)}\right] &= \int_{t_0}^{t_1} \nabla \cdot \mathbf{J_m(x,y,z,t)} \, dt \\
  - \nabla \cdot \mathbf{B(x,y,z,t_1)} + \nabla \cdot \mathbf{B(x,y,z,t_0)} &=  \int_{t_0}^{t_1} \nabla \cdot \mathbf{J_m(x,y,z,t)} \, dt \\
  - \nabla \cdot \mathbf{B(x,y,z,t_1)} &= - \nabla \cdot \mathbf{B(x,y,z,t_0)} + \int_{t_0}^{t_1} \nabla \cdot \mathbf{J_m(x,y,z,t)} \, dt \\
  - \nabla \cdot \mathbf{B(x,y,z,t_1)} &= - \left[\nabla \cdot \mathbf{B(x,y,z,t_0)} - \int_{t_0}^{t_1} \nabla \cdot \mathbf{J_m(x,y,z,t)} \, dt\right] \\
  \nabla \cdot \mathbf{B(x,y,z,t_1)} &= \nabla \cdot \mathbf{B(x,y,z,t_0)} - \int_{t_0}^{t_1} \nabla \cdot \mathbf{J_m(x,y,z,t)} \, dt \\
  \text{[Since \( \mathbf{J_m(x,y,z,t)} = 0 \) in nature]} \\
  \nabla \cdot \mathbf{B(x,y,z,t_1)} &= \nabla \cdot \mathbf{B(x,y,z,t_0)} - \int_{t_0}^{t_1} 0 \, dt \\
  \nabla \cdot \mathbf{B(x,y,z,t_1)} &= \nabla \cdot \mathbf{B(x,y,z,t_0)}
\end{align*}
\textbf{This proves that the divergence of the magnetic flux density function
is constant over time}. Since \( J_m(x,y,z,t_0) = 0 \) in nature, we can conclude
\( \frac{\partial}{\partial t} \rho_m(x,y,z,t_0) = 0 \) and subsequently
\( \nabla \cdot \mathbf{B(x,y,z,t_0)} = 0 \). This is \textbf{the origin of Gauss's Law
of Magnetism}.

\subsubsection{Ampere's Law}

Ampere's Law states that a time-varying electric function induces a magnetic
function. It is expressed as:
\begin{align*}
  \nabla \times \mathbf{H(x,y,z,t)} = \frac{\partial \mathbf{D(x,y,z,t)}}{\partial t} + \mathbf{J_e(x,y,z,t)} \tag{3}
\end{align*}
where
\begin{itemize}
  \item \( \mathbf{H(x,y,z,t)} \): magnetic function. \( \frac{A}{m} \)
  \item \( \mathbf{D(x,y,z,t)} \): electric flux density function.
    \( \frac{C}{m^2} \)
  \item \( \mathbf{J_e(x,y,z,t)} \): electric current density function.
    \( \frac{A}{m^2} \). In nature, \( \mathbf{J_e(x,y,z,t) \neq 0 } \) since
    electric monopoles that is, charges have been observed. It is defined as
    \( \mathbf{J_e(x,y,z,t)} = \rho_e(x,y,z,t) \mathbf{v(x,y,z,t)} \) where
    \( \rho_e(x,y,z,t) \) is the charge density function and
    \( \mathbf{v(x,y,z,t)} \) is the velocity function of the charge.
\end{itemize}
Let's have some fun with this equation.
\begin{align*}
  \text{[Divergence of both sides]} \\
  \nabla \cdot (\nabla \times \mathbf{H(x,y,z,t)}) &= \nabla \cdot \left(\frac{\partial \mathbf{D(x,y,z,t)}}{\partial t} + \mathbf{J_e(x,y,z,t)}\right) \\
  \nabla \cdot (\nabla \times \mathbf{H(x,y,z,t)}) &= \nabla \cdot \frac{\partial \mathbf{D(x,y,z,t)}}{\partial t} + \nabla \cdot \mathbf{J_e(x,y,z,t)} \\
  \nabla \cdot (\nabla \times \mathbf{H(x,y,z,t)}) &= \frac{\partial}{\partial t} (\nabla \cdot \mathbf{D(x,y,z,t)}) + \nabla \cdot \mathbf{J_e(x,y,z,t)} \\
  0 &= \frac{\partial}{\partial t} (\nabla \cdot \mathbf{D(x,y,z,t)}) + \nabla \cdot \mathbf{J_e(x,y,z,t)} \\
  - \frac{\partial}{\partial t} (\nabla \cdot \mathbf{D(x,y,z,t)}) &= \nabla \cdot \mathbf{J_e(x,y,z,t)} \\
  \text{[Integrating both sides over time]} \\
  \int_{t_0}^{t_1} - \frac{\partial}{\partial t} (\nabla \cdot \mathbf{D(x,y,z,t)}) \, dt &= \int_{t_0}^{t_1} \nabla \cdot \mathbf{J_e(x,y,z,t)} \, dt \\
  - \left[\nabla \cdot \mathbf{D(x,y,z,t)}\right]_{t_0}^{t_1} &= \int_{t_0}^{t_1} \nabla \cdot \mathbf{J_e(x,y,z,t)} \, dt \\
  - \left[\nabla \cdot \mathbf{D(x,y,z,t_1)} - \nabla \cdot \mathbf{D(x,y,z,t_0)}\right] &= \int_{t_0}^{t_1} \nabla \cdot \mathbf{J_e(x,y,z,t)} \, dt \\
  - \nabla \cdot \mathbf{D(x,y,z,t_1)} + \nabla \cdot \mathbf{D(x,y,z,t_0)} &=  \int_{t_0}^{t_1} \nabla \cdot \mathbf{J_e(x,y,z,t)} \, dt \\
  - \nabla \cdot \mathbf{D(x,y,z,t_1)} &= - \nabla \cdot \mathbf{D(x,y,z,t_0)} + \int_{t_0}^{t_1} \nabla \cdot \mathbf{J_e(x,y,z,t)} \, dt \\
  - \nabla \cdot \mathbf{D(x,y,z,t_1)} &= - \left[\nabla \cdot \mathbf{D(x,y,z,t_0)} - \int_{t_0}^{t_1} \nabla \cdot \mathbf{J_e(x,y,z,t)} \, dt \right] \\
  \nabla \cdot \mathbf{D(x,y,z,t_1)} &= \nabla \cdot \mathbf{D(x,y,z,t_0)} - \int_{t_0}^{t_1} \nabla \cdot \mathbf{J_e(x,y,z,t)} \, dt \\
\end{align*}
\textbf{This proves that the divergence of the electric flux density function
changes over time}. To find anything meaningful we must find more information
about the divergence of the electric current density function. This is given
to us by the \textbf{charge continuity equation}:
\begin{align*}
  Q(t) &= \int_V \rho_e(x,y,z,t) \, dV \hspace{1cm} \text{[Total Charge in a Volume]} \\
  \text{[Rate of Change of Total Charge]} \\
  \frac{dQ(t)}{dt} &= \int_V \frac{\partial \rho_e(x,y,z,t)}{\partial t} \, dV \hspace{1cm} \text{[In a Volume]} \\
  \frac{dQ(t)}{dt} &= - \oint_S \mathbf{J_e(x,y,z,t)} \cdot d\mathbf{S} \hspace{1cm} \text{[Out of a Volume]} \\
  \text{[Continuity; In = Out]} \\
  - \oint_S \mathbf{J_e(x,y,z,t)} \cdot d\mathbf{S} &= \int_V \frac{\partial \rho_e(x,y,z,t)}{\partial t} \, dV \\
  - \int_V \nabla \cdot \mathbf{J_e(x,y,z,t)} \, dV &= \int_V \frac{\partial \rho_e(x,y,z,t)}{\partial t} \, dV \hspace{1cm} \left[\oint_S \mathbf{J_e} \cdot d\mathbf{S} = \int_V \nabla \cdot \mathbf{J_e} \, dV\right] \\
  \text{[Differentiating both sides with respect to Volume]} \\
  \frac{\partial}{\partial V} \left(- \int_V \nabla \cdot \mathbf{J_e(x,y,z,t)} \, dV\right) &= \frac{\partial}{\partial V} \left(\int_V \frac{\partial \rho_e(x,y,z,t)}{\partial t} \, dV\right) \\
  - \nabla \cdot \mathbf{J_e(x,y,z,t)} &= \frac{\partial \rho_e(x,y,z,t)}{\partial t} \tag{4}
\end{align*}
Substituting (4) into (3), we get
\begin{align*}
  \nabla \cdot \mathbf{D(x,y,z,t_1)} &= \nabla \cdot \mathbf{D(x,y,z,t_0)} - \int_{t_0}^{t_1} \nabla \cdot \mathbf{J_e(x,y,z,t)} \, dt \\
  \nabla \cdot \mathbf{D(x,y,z,t_1)} &= \nabla \cdot \mathbf{D(x,y,z,t_0)} - \int_{t_0}^{t_1} - \frac{\partial \rho_e(x,y,z,t)}{\partial t} \, dt \\
  \nabla \cdot \mathbf{D(x,y,z,t_1)} &= \nabla \cdot \mathbf{D(x,y,z,t_0)} + [\rho_e(x,y,z,t_1) - \rho_e(x,y,z,t_0)] \\
  \nabla \cdot \mathbf{D(x,y,z,t_1)} &= \nabla \cdot \mathbf{D(x,y,z,t_0)} + \rho_e(x,y,z,t_1) - \rho_e(x,y,z,t_0)
\end{align*}
We find that the divergence of the electric flux density function changes as
the charge density function changes over time \( \nabla \cdot \mathbf{D(x,y,z,t)} =
\rho_e(x,y,z,t) \). This is \textbf{the origin of Gauss's Law of Electricity}.

\subsection{Static and Quasi-Static}

Faraday's and Ampere's Law (alongisde Gauss) fully define any electromagnetic system
at all times and space. Both equations are \textbf{first order partial differential
equations (PDEs)} which \textbf{are coupled with each other}. Sadly, for solving the
full system is mathematically very expensive. One therefore wonders, \textbf{Is it
possible to simplify and dumb down?}


\textbf{Yes, but only under conditions}. This is where Static and Quasi-Static
come into play. We need to make some assumptions.
\begin{enumerate}
  \item \textbf{Static}: We assume that all functions are frozen in time. This is the
    strongest assumption we can make and it reduces our PDEs to \textbf{ordinary
    differential equations (ODEs)} in space only. This is the DC regime where all functions
    are time-independent. To achieve this, the rate of change of all functions must
    be zero:
    \begin{align*}
      \frac{\partial \mathbf{B(x,y,z,t)}}{\partial t} = 0 \\
      \frac{\partial \mathbf{D(x,y,z,t)}}{\partial t} = 0 \\
      \frac{\partial \mathbf{E(x,y,z,t)}}{\partial t} = 0 \\
      \frac{\partial \mathbf{H(x,y,z,t)}}{\partial t} = 0 \\
      \frac{\partial \mathbf{J_e(x,y,z,t)}}{\partial t} = 0 \\
      \frac{\partial \mathbf{J_m(x,y,z,t)}}{\partial t} = 0
    \end{align*}
  \item \textbf{Quasi-Static}: We assume that the flux functions [\(\mathbf{B}\) \&
    \(\mathbf{D}\)] do not change with time but we allow other functions [\(\mathbf{E}\),
    \(\mathbf{H}\), \(\mathbf{J}_e\), and \(\mathbf{J}_m\)] to change with time. This is
    a weaker assumption and it reduces our PDEs to ODEs in space only for the flux
    functions. This is the quasi-static regime. To achieve this, we can set the time
    derivatives of the flux functions to zero:
    \begin{align*}
      \frac{\partial \mathbf{B(x,y,z,t)}}{\partial t} = 0 \\
      \frac{\partial \mathbf{D(x,y,z,t)}}{\partial t} = 0
    \end{align*}
\end{enumerate}
In either case, we can simplify Faraday's and Ampere's Law to:
\begin{align*}
    -\nabla \times \mathbf{E(x,y,z,t)} &= \frac{\partial \mathbf{B(x,y,z,t)}}{\partial t} + \mathbf{J_m(x,y,z,t)} \hspace{1cm} \text{[Faraday's Law]} \\
                                       &= 0 + \mathbf{J_m(x,y,z,t)} \hspace{1cm} \left[\frac{\partial \mathbf{B(x,y,z,t)}}{\partial t} = 0\right] \\
                                       &= \mathbf{J_m(x,y,z,t)} \\
     &= 0  \hspace{1cm} \text{[since \( \mathbf{J_m(x,y,z,t)} = 0 \)]} \tag{5} \\\\
    \nabla \times \mathbf{H(x,y,z,t)} &= \frac{\partial \mathbf{D(x,y,z,t)}}{\partial t} + \mathbf{J_e(x,y,z,t)} \hspace{1cm} \text{[Ampere's Law]} \\
                                      &= 0 + \mathbf{J_e(x,y,z,t)} \hspace{1cm} \left[\frac{\partial \mathbf{D(x,y,z,t)}}{\partial t} = 0\right] \\
                                      &= \mathbf{J_e(x,y,z,t)} \tag{6}
\end{align*}

Equations (5) and (6) are \textbf{instantaneous spatial statements}. They hold at every
value of \(t\) individually but say nothing about how \(\mathbf{E}\) or \(\mathbf{H}\)
behave \emph{between} t's. A function can satisfy \(\nabla\times\mathbf{E}=0\) at
every instant and still change with time. For example \(\mathbf{E} = \sin(\omega t)\,
\nabla\phi(x,y,z)\) has zero curl for all \(t\), yet clearly changes with time. So (5)
and (6) do not, by themselves, let us drop \(t\) from \(\mathbf{E}\), \(\mathbf{H}\),
or \(\mathbf{J}_e\) (\& \(\mathbf{J}_m\)).

\begin{enumerate}
  \item \textbf{Equation 5}: At every instant (t), the electric function  has zero
    curl = it is irrotational = it is conservative = the potential difference between
    two points does not depend on the path taken between them = total work done by
    the electric function on a charge moving from point a to point b and back to
    point a is zero, \emph{at that instant}.
    \begin{align*}
      \text{[At every instant \(t\)]} \\
      -\nabla \times \mathbf{E(x,y,z,t)} &= 0 \\
      \nabla \times (\nabla V(x,y,z,t)) &= 0 \hspace{1cm} \left[\nabla \times (\nabla F) = 0 \right] \\
      -\nabla \times \mathbf{E(x,y,z,t)} &= \nabla \times (\nabla V(x,y,z,t)) \\
      \mathbf{E(x,y,z,t)} &= -\nabla V(x,y,z,t) \\
      \text{[Line Integral of Electric Function A to B to A]} \\
      \oint_a^b \mathbf{E(x,y,z,t)} \cdot d\mathbf{l} &= \oint_a^b -\nabla V(x,y,z,t) \cdot d\mathbf{l} \\
                                                &= \int_a^b -\nabla V(x,y,z,t) \cdot d\mathbf{l} + \int_b^a -\nabla V(x,y,z,t) \cdot d\mathbf{l} \\
                                                &= -\left[ V(b,t) - V(a,t) \right] - \left[ V(a,t) - V(b,t) \right] \\
                                                &= 0
    \end{align*}
    This is the origin of Kirchhoff's Voltage Law (KVL): the total potential difference
    around any closed loop is zero. Notice this did not require \(\mathbf{E}\) or \(V\)
    to be time-independent, only that \(\nabla\times\mathbf{E}=0\) holds instant by
    instant. This is precisely why KVL remains valid in the quasi-static (slowly
    time-varying, \(L \lll \lambda\)) regime, not just in true DC.

  \item \textbf{Equation 6}: Taking the divergence of both sides of
    equation (9), at each instant \(t\):
    \begin{align*}
      \nabla \cdot (\nabla \times \mathbf{H(x,y,z,t)}) &= \nabla \cdot \mathbf{J_e(x,y,z,t)} \\
      0 &= \nabla \cdot \mathbf{J_e(x,y,z,t)} \hspace{1cm} \text{[since divergence of a curl is always zero]} \tag{11}
    \end{align*}

    \begin{itemize}
      \item \textbf{Charge Density Is Genuinely Time-Independent}:
        Substituting (5) into (11):
        \begin{align*}
          0 &= -\frac{\partial \rho_e(x,y,z,t)}{\partial t} \tag{12}
        \end{align*}
        Unlike \(\mathbf{E}\), \(\mathbf{H}\), and \(\mathbf{J}_e\)
        above, this is a genuine, non-instantaneous result: it says
        \(\rho_e\)'s rate of change is zero \emph{for all} \(t\), not
        just at a single instant. This is the one function we can
        honestly write without \(t\):
        \begin{align*}
          \rho_e(x,y,z,t) = \rho_e(x,y,z) \tag{12a}
        \end{align*}

      \item \textbf{Zero Total Current Out of a Volume}: Integrating
        both sides of equation (11) over a volume \(V\), at a fixed
        instant \(t\):
        \begin{align*}
          \int_V \nabla \cdot \mathbf{J_e(x,y,z,t)} \, dV &= \int_V 0 \, dV \\
          \int_V \nabla \cdot \mathbf{J_e(x,y,z,t)} \, dV &= 0 \\
          \oint_S \mathbf{J_e(x,y,z,t)} \cdot d\mathbf{S} &= 0 \hspace{1cm} \text{[by Divergence Theorem]} \tag{13}
        \end{align*}
    \end{itemize}
    Equation (13) shows that the total current flowing out of a
    closed surface is zero \emph{at every instant \(t\)} -- again,
    without requiring \(\mathbf{J}_e\) itself to be constant in
    time. Equation (12) shows the total charge enclosed by that
    surface genuinely does not change with time. Together, these
    mean the current flowing into a volume equals the current
    flowing out, at every instant. This is the origin of Kirchhoff's
    Current Law (KCL).
\end{enumerate}

\subsubsection{E, D and H, B Functions}

Electric Flux Density \( \mathbf{D(x,y,z,t)} \) and Electric Field \( \mathbf{E(x,y,z,t)} \)
are related to each other by the permittivity of the medium \( \epsilon \)
and the polarization of the medium \( \mathbf{P} \):
\begin{align*}
  \mathbf{D(x,y,z,t)} = \epsilon \mathbf{E(x,y,z,t)} + \mathbf{P}
\end{align*}

Magnetic Flux Density \( \mathbf{B(x,y,z,t)} \) and Magnetic Field \( \mathbf{H(x,y,z,t)} \)
are related to each other by the permeability of the medium \( \mu \)
and the magnetization of the medium \( \mathbf{M} \):
\begin{align*}
  \mathbf{B(x,y,z,t)} = \mu \mathbf{H(x,y,z,t)} + \mathbf{M}
\end{align*}
